\section{Lean formalization and object-theory endpoints}
\label{sec:lean-route}

\subsection{Two independently usable representations}
The file \texttt{lean/FiniteInversion.lean} uses only Lean's bundled
library. It was checked with Lean 4.32.0, the version pinned by ProveIt,
and with Lean 4.19.0. Entries of its matrix are genuine Lean integers.
If $W(i,j)$ is the weight at row $i$, column $j$, write
\[
 (Lf)(i)=\sum_{j<i}W(i,j)f(j),\qquad (Zf)(i)=f(i)+(Lf)(i).
\]
Define $b_0(i)=0$, $b_{t+1}(i)=g(i)-(Lb_t)(i)$, and $b(i)=b_{i+1}(i)$.
If $i<t$, induction on $t$ gives $b_{t+1}(i)=b_t(i)$, since every
previous entry needed in row $i$ has index $j<i$. Thus $Zb=g$, and
substitution, from the lowest row upwards, gives $b=f$ whenever $g=Zf$.
These are \texttt{integrate\_invert} and \texttt{invert\_integrate}.
The file also proves uniqueness on a specified prefix of rows.

Subtract two consecutive iterates. The first difference is $g$, and
each remaining difference is $-L$ applied to its predecessor. Hence
\[
 b(i)=\sum_{r=0}^{i} ((-L)^r g)(i).
\]
Each summand is a signed sum over strictly increasing length-$r$ chains,
by expanding its last edge. For $r>i$ the row vanishes.
The file checks both assertions uniformly in $i,r,W,g$.
With $W(i,j)=[j+1\mid i+1]$ and $g(i)=[i=0]$, the recurrence computes
$\mu(i+1)$. Its identification with the prime-factor definition follows
from the divisor cancellation proof and uniqueness. The Lean file proves
divisor cancellation and the chain formula for its recurrence-defined
function. A second proof, \texttt{ExponentBox.box\_cancellation}, sums the
prime-factor definition over a list of exponent ranges of arbitrary
length. The coordinate factor at $e=0,1,\ge2$ is respectively $1,-1,0$.
Its prefix sum up to exponent $a$ is $[a=0]$. Factoring this sum off
the first coordinate and iterating proves the product formula. The human
divisor--exponent bijection connects this calculation with the numerical
divisor scan; the standalone file leaves that bijection to the report.

The companion \texttt{FormalSeries.lean} imports this file. The coefficient
calculations in Section~\ref{sec:mob-basel} describe its scaled
representations and conclusions. Run \texttt{sh lean/check.sh} to check
both files together.

\subsection{The existing internal proofs}
For object proofs use revision \texttt{65b76a5b8aa8} of ProveIt and the
namespace \texttt{LeanProofs.}\allowbreak\texttt{IPRA.}\allowbreak\texttt{AnalyticNumberTheory}. The relevant
source directory is
\href{https://github.com/VladimirReshetnikov/ProveIt/tree/65b76a5b8aa87ecad4103c2d287eff5660d3c84c/Logic/PrimitiveRecursiveArithmetic/Derivations/Lean/IPRADerivations/AnalyticNumberTheory}{the analytic number theory derivations directory}.
The ordinary calculus has function codes \texttt{PR k}, terms
\texttt{PTm n} over $n$ variables, formulas \texttt{Fm PTm n}, and
deductions \texttt{Prf PRAAx} indexed by their contexts and conclusions.
A term body with free variables for its input and parameters expresses a
schematic arithmetic function. Object theorems use numerical variables
for bounds. They may be specialized at arbitrary terms in any context
which supplies the stated hypotheses.

The repository takes $\mu_N^+(x)=[s_N(x)=0][\omega_N(x)\text{ even}]$ and
$\mu_N^-(x)=[s_N(x)=0][\omega_N(x)\text{ odd}]$, scanning primes $q<N$.
Here $\omega_N$ counts nonzero valuations, and $s_N$ counts those at
least two. Once $N>x\ge1$ the pair is stable. Integer equations are
written by transferring negative components across the equality.
The following are source-inspected object-theory declarations at that
revision. The replay file lists their fully resolved module imports.
\begin{itemize}
\item \texttt{Mobius.swap}, \texttt{Mobius.kill},
\texttt{Mobius.mobius\_one}, and \texttt{Mobius.mobius\_cancel} give
the prime steps and the divisor sum. \texttt{MobiusBounds.cancel\_gen}
allows an arbitrary prime bound and an arbitrary divisor scan limit,
both strictly above the positive argument. This matters when one cutoff
is shared by many terms.
\item \texttt{MobiusInversion.mobius\_inversion} instantiates the divisor
inversion at a term body and \texttt{phi\_mobius} derives its totient
specialization. For signed input bodies use the theorem on each natural
component and add the cross terms.
\item \texttt{SummatoryPhi.summatory\_phi} gives the floor-square
totient identity. \texttt{MuLambda.mu\_lambda} gives
$\mu*\Lambda_w=-\mu L_w$ for the prime weight represented by the
given body. \texttt{PsiModulus.}\allowbreak\texttt{psi\_modulus} supplies the quantitative
PNT for the repository's finite harmonic weights; \texttt{FinalPNT.}\allowbreak\texttt{arith\_pnt}
gives its prime-counting version.
\end{itemize}
These existing derivations and the new standalone Lean proofs address
different judgments. We inspected the pinned upstream sources; this
packet includes replay commands for a prepared checkout. We ran both
standalone files locally. The remaining object derivations described below
are a formalization plan.

\subsection{Constructing additional internal deductions}
\begin{enumerate}
\item \emph{Encoding and induction.} Use \texttt{PR.comp} to compose
body codes and \texttt{PR.prec} with arguments in the order
(counter, parameters), its step receiving (predecessor, previous value,
parameters). Use bounded sums and products from \texttt{Core/Definitions}
and \texttt{Core/MultiSums}, \texttt{Core/Search} for the first matching
index, and \texttt{Core/Fubini} for two-index rearrangements. Define
equality of signed pairs by $u_++v_-=v_++u_-$. For a bounded family of
integer equalities $E_i$, form its failure count
$\sum_{i<t}(1-[E_i])$. Induct on its equality to zero using
\texttt{PRA.qfInd}. Nested counts allow several bounded indices.
Recover each equality since a nonzero summand contributes at least one.
This implements every use of finite strong induction on the prefix.
\item \emph{Inversion and posets.} Fix a packed finite relation and a
numerical bound $B$. On $i<B$ implement $L$ by the displayed bounded
sum. Store the $B$ entries of the current row vector as a finite list
of signed pairs; compute the next vector by one length-$B$ fold.
The computation of $b_{i+1}(i)$ now recurses on a natural code at fixed
$B$ instead of iterating over functions. Induct on the zero failure count
of the equalities $b_{t+1}(j)=b_t(j)$ for $j<\min(B,t)$.
Base, step and uniqueness are the finite proofs above. For an interval
$[a,b]$ in a poset use a finite list with $x<y$ implying earlier
position for $x$. Set $W(i,j)=[x_j<x_i]$ and $g(i)=[x_i=a]$.
On divisibility take $B=n$, $x_i=i+1$, and instantiate
\texttt{MobiusBounds.cancel\_gen} to compare the two definitions of
$\mu$ pointwise by prefix uniqueness.
\item \emph{Coefficient identities.} Store degrees $0$ through $B$,
representing a polynomial by a signed coefficient list. Convolution at
$k\le B$ is the fold $\sum_{j\le k}a_j b_{k-j}$. For Dirichlet
products replace the additive incidence by $[de=k]$ and scan only
$1\le d,e\le B$. For the Lambert series the coefficient of $X^k$
adds one term per $d$ with $d\mid k$. Apply
\texttt{MobiusBounds.cancel\_gen} at that $k$ and close all degrees
through $B$ by the zero failure count. Formal logarithms may be stored
with common denominator $(B+1)!$; prove divisibility before exact
division. Polynomial cyclotomic arguments use lists with the explicitly
given degree bound, monic division with remainder, and equality of
coefficients. Finite geometric sums are first proved by recursion on
their length, so evaluating at a unit in a specified quotient ring
uses only addition and multiplication of integer pairs or lists.
For the irreducible count over a field of cardinality $q$, supply two
$q$-by-$q$ operation tables and validate the field axioms by a bounded
scan. Represent a monic of degree $n$ by the base-$q$ digits of a number
less than $q^n$. For $n>0$ its irreducibility bit tests every monic
divisor of degree $1$ through $\lfloor n/2\rfloor$ by remainder.
Summing this bit gives $I_q(n)$. If a division splits a polynomial of
positive degree, both new degrees decrease. Recursion through a coded
list, with sum of the pending degrees bounded by $n$, gives a complete
factor list in at most $2n-1$ node visits. Use the Euclidean B\'ezout
identity to prove Euclid's lemma for monic irreducibles, then compare
multiplicities. This supplies the factorization bijection used by the
finite Euler product.
\item \emph{Finite fields and Fourier sums.} Encode the two $q$-by-$q$
field operations by lists of entries $<q$. Test associativity, distributivity
and the inverse laws on the bounded lists of tuples. Division by a monic
polynomial is a fold on its degree, stripping the leading term. Record
quotient and remainder and verify the product identity at each step.
For degree $d$ enumerate coefficient codes $<q^d$ and enumerate the
possible monic factors through degree $\lfloor d/2\rfloor$; the bit
rejecting all of them defines $I_q(d)$. Induct on degree to establish
the unique factor multiset, then count products in the coefficient table
through degree $n$ and apply the existing inversion theorem to the term
$d I_q(d)$. For $C_n(h)$, supply the Euclidean-algorithm certificate
$(U,V)$ for each invertible difference. In the quotient ring the checker
verifies its inverse by remainder equality. The same geometric-sum proof
over the table of residues modulo $p$ gives the congruence.
\item \emph{Approximation statements.} A rational name has primitive
recursive numerator and positive denominator at each requested precision.
Implement a complex name by two such pairs. For a fixed input code or a
specified PR family and precision $k$, substitute the displayed cutoff in
the relevant finite sum or product, perform exact denominator clearing,
and express the error as two integer inequalities per real coordinate.
For a tail use an additional free end index $V\ge U$, and prove the
displayed estimate for every such $V$ by induction on $V-U$. Pass
between different cutoffs by the triangle inequality and the given
moduli. At the PNT step instantiate \texttt{psi\_modulus} with the
repository's own mesh and denominator; compare alternative rational
meshes by refining to their product. These operations leave all
induction formulas quantifier-free after the bounded checks are compiled.
A statement about a family of names quantifies over indices and precisions
for that fixed family.
\item \emph{The strong summatory estimate and boundary moments.}
The ordinary PNT modulus suffices for Section~\ref{sec:pnt}.
For Theorem~\ref{thm:strongM} the stronger analytic input in \cite{strip}
must itself be represented by finite deductions. Implement the Euler
summation samples with rational logarithm and exponential enclosures,
the fourth-power strip comparison, and the regular reciprocal evaluation
at the pole. The finite contour invariant equates the signed internal
edge sum to zero, then accumulates at most one Taylor error per cell.
Apply it to the smoothed signed sum at the two rational arguments.
After raising the unsmoothing estimate to an eighth power, the closed
matrix compares $2^k|M(x)|^8$ and $x^8$. To remove arbitrary precision
from an inequality between integers it suffices to make the additive
enclosure slack less than $1$: discreteness then gives the exact weak
inequality. The moment target is the rational partial-sum comparison
following \eqref{eq:momentTail}. A second possible dependency is the
strong $\vartheta$ modulus and the bounded numerical induction in
Section~\ref{sec:bootstrap}. Both routes require additional embedded
deductions.
\end{enumerate}
